\documentclass[acmsmall,screen,review,anonymous,nonacm]{acmart}
\usepackage{silence}
\WarningFilter{latex}{Command \showhyphens has changed}


\title{how to run knitting machine}
\author{TODO}

\begin{document}

\maketitle

\section{Helpful links}
\begin{itemize}
    \item \href{https://fablab.cs.washington.edu/}{fablab site}
    \item \href{../SOP.pdf}{Knitting machine SOP (this doc is basically a port of it)}
    \item \href{https://wiki.computationalmachineknitting.org/app/page/1GeyC58yhlU_17IHSqfhCyQipoM_dPPBK?p=1GeyC58yhlU_17IHSqfhCyQipoM_dPPBK}{Computational knitting machine wiki}
\end{itemize}

\section{Setup}
\label{sec:setup}
\begin{itemize}
    \item Make sure Nat is in the room for now. We can figure out how to get Elden trained to
        run the machine solo later
    \item Make sure the air is on, and the machine is plugged in
    \item Carriage should be on the left and selector pins should all be down
    \item Yarn tensioning device should be twist-free and have nothing dangling into the
        machine; yarns shouldn't be tangled with each other (somebody keeps tangling the yarns!!!)
    \item Covers should be closed
    \item Power on: \begin{enumerate}
            \item Twist the big red knob
            \item After it boots, press the clear ``drive power'' switch
            \item If it asks ``DRIVE WITH OPERATION SWITCH'', hit the green ``operation switch''. It
                should zero out the motors
            \item Hit the pink ``Prepa'' button, then ``reset'' (done by pressing ``F4'')
        \end{enumerate}
    \item If it's the first machine run of the day (safe to assume it is), do debugging then needle training: \begin{enumerate}
            \item Press maintenance (driver + wrench on touchscreen)
            \item Debug mode
            \item Highlight test pattern (``TEST PAT.'')
            \item Press go (``F5'')
            \item Let it run for at least 10 seconds (makes sure carriage is in a valid state), stop with ``F2''. Debug done!
            \item Needle training: tap needle training in the maintenance menu
            \item Tap start
            \item Green button
        \end{enumerate}
\end{itemize}

\section{DAT to 000}

\begin{itemize}
    \item Get a FAT32 flash stick. Usually there's one lying around the lab or in the knitting machine. Nat also has one
    \item Put a ``.DAT'' file onto the flash stick. ``.DAT'' is Shima Seiki's assembly language; knitout has a compiler
        to ``.DAT''.
    \item Plug the flash stick into the APEX computer; double click the ``.DAT'' to open it in KnitPaint.
        You'll need to press ``Exec'' to accept the defaults tfor how ot open it
    \item Click the ``Auto'' button at the top of the screen. Make sure there's a dotted line around
        the instructions, NOT the options (not the colored columns on either side). If not, click the instructions to select it.
        Click OK
    \item If you want to change carriers:
        \begin{itemize}
            \item Don't use the ``Carrier Change'' button
            \item Manually change the stuff in the ``Yarn'' tab in the ``N'' row
            \item Do the same changes for ``Y. Hold Hook''
            \item The ``YF'' should be magically updated too now
        \end{itemize}
    \item Make sure the output device and file name is to your likingg (output to flash stick)
    \item Click ``Execute'' to save to ``.000''. Accept/acknowledge any warnings here
    \item It should come up with an ``Output Data'' window. Click ``OK'', not ``Skip''
    \item It should have a ``Knit Simulation'' window. Always click ``Knitting Assist'' to have it look
        for problems. No assist information is good
    \item You can click the play button after Knitting Assist to watch the simulation go. You can use
        the ``Area'' button (icon in top left of ``KnittingViewer'' window) to select a part of the
        knit instructions that you want to focus in on
    \item Once you're happy, eject the USB with the ``Remove USB'' button on the top right of the
        main screen and yank the USB out
\end{itemize}

\section{Running a 000}

\begin{itemize}
    \item Plug the FAT32 flash stick with the ``.000'' file into the machine in its fancy USB port (machine can be on or off)
    \item Run machine setup~\ref{sec:setup} if you haven't already
    \item Tab the USB icon -- the menu should look a little glitched out
    \item Tap the ``upload from folder'' looking icon
    \item Tap the ``.000'' file that you want, then press ``F5'' and ``EXEC''
    \item Filename should be listed on the homepage
    \item Press ``Prepa'', then ``F5'' to prepare to run. It should beep when it's done
    \item Green to go
    \item Red to pause, green to resume, white is ``ultra low'' speed
\end{itemize}

\section{Turning off the machine}

\begin{itemize}
    \item Hold the clear power button until it beeps and says it powered off
    \item Turn the big red knob
\end{itemize}

\section{Knit object is stuck in there!}

Need to press off. Find Nat.

\section{Something wacky is happening!}

Find Nat.

\bibliographystyle{ACM-Reference-Format}
\bibliography{references}

\end{document}
